\subsection*{Question 1}
	$100 - 80  =  20$.
	\quad
	La réduction est de 20\,\euro{}.

	Le pourcentage de réduction est donc :\quad
		\begin{tabular}{c|c|c}
	&en \euro  &en ~\% \\\hline
		réduction	& 20 &$x$  \\\hline
		prix initial&100&100\\
	\end{tabular}\qquad
	\mbox{$x = \dfrac{20\times100}{100} = 20$},

	\textbf{Bonne réponse : B.}

\subsection*{Question 2}
	10~\% de 60~\%, c'est 6~\% du total, donc la proportion est \np{0.06}.

	\textbf{Bonne réponse : C.}

\subsection*{Question 3}
	Coefficients multiplicateurs : \quad $1 + \dfrac{10}{100}  =  1,1$ : hausse de 10\,\%\quad $1 - \dfrac{20}{100}  = 0,8$ : baisse de 20\,\%.

	Le prix a donc été multiplié par \np{1.1}, puis par \np{0.8}, donc par le coefficient global : \quad $c_g  =   1,1 \times 0,8  =  \np{0.88}$.

	$t_g  =  c_g - 1  =  0,88 - 1  =  -0,12  =  -\dfrac{12}{100}$.

	Cela correspond à une baisse de 12~\%.

	\textbf{Bonne réponse : B.}

\subsection*{Question 4}
	La température d'ébullition en ~\textdegree{}F est donc \quad  \mbox{$\np{1.8}\times100+32 = 180+32 = 212$}.

	\textbf{Bonne réponse : D.}

\subsection*{Question 5}
	On sait que \quad \mbox{$F = \np{1.8} C+32$},\quad  donc  \quad \mbox{$F-32 = \np{1.8} C$} \quad et \quad \mbox{$C = \dfrac{F-32}{\np{1.8}}$}.

	\textbf{Bonne réponse : B.}

\subsection*{Question 6}
	Calculons :\quad \mbox{$A = \dfrac{\dfrac53}{\dfrac{15}{1}} = \dfrac53\times\dfrac{1}{15} = \dfrac{5\times1}{3\times3\times5} = \dfrac19$}.

	\textbf{Bonne réponse : A.}

\subsection*{Question 7}
	La droite n'est pas parallèle à l'axe des ordonnées, elle a donc une équation réduite du type \mbox{$y = ax + b$}.

	Elle passe par les points \mbox{$(0;2)$} et \mbox{$(4;0)$}, donc sa pente est \mbox{$a = \dfrac{0-2}{4-0} = -\np{0.5}$} et son ordonnée à l'origine \mbox{$b = 2$}.

	\textbf{Bonne réponse : B.}

\subsection*{Question 8}
	Le coefficient directeur de (EF) est \mbox{$\dfrac{y_\text{E}-y_\text{F}}{x_\text{E}-x_\text{F}} = \dfrac{25-15}{20-5} = \dfrac{10}{15} = \dfrac23$}.

	\textbf{Bonne réponse : C.}

\subsection*{Question 9}
	La fonction est positive sur \mbox{$[-5;-1]$}, elle change de signe en $-1$ et est négative sur \mbox{$[-1;3]$}.

	\textbf{Bonne réponse : B.}

\subsection*{Question 10}
	Le troisième quartile d'une série est la plus petite valeur telle que les \mbox{$\dfrac34$} au moins des valeurs de la série sont inférieures ou égales à celle-ci.

	Donc la proportion d'élèves ayant une note inférieure ou égale à 12 est d'au moins 75~\%.

	\textbf{Bonne réponse : C.}

\subsection*{Question 11}
	La proportion de joueurs de moins de 24 ans est égale (en pourcentage) à \mbox{$\np{8.5}+\np{12.5}+\np{18.9} = \np{39.9}$}, soit \np[\%]{39.9}

	\textbf{Bonne réponse : C.}

\subsection*{Question 12}
	On choisit un élève au hasard parmi les garçons, donc on est en situation d'équiprobabilité : les proportions sont assimilables à des probabilités.

	Il y a 500 garçons, dont 200 ne pratiquent aucune activité sportive, donc la probabilité recherchée est \mbox{$\dfrac{200}{500}$}.

	\textbf{Bonne réponse : D.}
